\documentclass{memoir}

\title{Dynamic Distributed Minimum Spanning Tree}

\author{Prof. Gabriele Puppis\\ DMIF, University of Udine, Italy}

\date{Version 0.3, \today}

\begin{document}


%\begin{titlingpage}
\maketitle
\begin{abstract}
The aim of this document is to describe how to write a report for the project-exam for the course Distributed Systems, at the University of Udine.
It is also a guideline for the development of the project itself.

This document is very brief and succint, and it is by no means comprehensive of all the informations that should be given about a software project. You are invited to supplement these guidelines as needed to best describe your work.
\end{abstract}
%\end{titlingpage}

\chapter{Introduction}\label{ch:intro}

\section{Description of the problem}
% A description of the application.

This project addresses the problem of maintaining a Minimum-Weight Spanning Tree (MST) in
a distributed network with dynamic topology. Specifically:

\begin{itemize}
\item links can fail and recover at arbitrary times, possibly concurrently,
\item no central coordinator exists: each node has knowledge only of its local
environment.
\end{itemize}

The goal is that, after any sequence of topological changes, the system converges to a
correct MST of the current network within finite time.

\section{Structure of the implementation}
% The overall structure of the implementation: how resources are deployed, which are the players, the roles.

\section{Features and Algorithms}
% The distributed system features (and the transparencies) and algorithms you intend to implement.

Features:
- fault tolerance
- asynchrony
- ...

\section{Testing}
% Your plan for testing the system.

\section{Schedule}
% A schedule for how you plan to carry our your design and implementation

\section{Glossary}

\begin{itemize}
\item Node
\item Link
\item Fragment
\end{itemize}


\chapter{Analysis}\label{ch:analysis}

In this chapter, we describe in detail functional and non-functional requirements of a solution for the problem.

\section{Functional requirements}
Which functions must be offered to users / other programs?  Which are the input data and the output data? Which is the expected effect? 

\section{Non functional requirements}

Cheng et al. : PA:
- eventual consistency: the system reaches a consistent state after a finite time
- availability (even if possibly not optimal)
- partition tolerance: non-negotiable since we want a dynamic mst

Everything about mode and transparencies: availability, mobility, security, fault tolerance, etc.

Are there execution time bounds? Minimum data rates?

If requested, specific platforms/languages/middlewares requirements for the implementation can be decided here. (E.g.: if the project is on a SOA, we may request that functions are offered via SOAP or RESTful services). 



\chapter{Project}

This chapter is devoted to the description of the general architectures, and specific algorithms.

\section{Logical architecture}
Describe the components of your systems: modules/objects/components/services.
For each component, describe the functionalities it implements, and by who is used.

\subsection{Protocols and Algorithms}
% Communication between components.  UML sequence diagrams go here.
% Also, put here a detailed description of distributed algorithms used to solve specific problems of the project.

\subsubsection{Messages}
Table of all exchanged messages (maybe this it is better to insert this section in chap 3)

% -------------------------------------------------------
\subsubsection{Failure response protocol}

% First of all let's consider a network which have already converged to a MST. When a link
% fails:
% - the MST is broken into two fragments,
% - the two nodes connected by that link will detect the failure,
% - a protocol to find a new edge to connect the two fragments will be executed
% . The protocol is based on the GHS algorithm, which is a distributed algorithm for finding a minimum spanning tree in a connected, undirected graph. The GHS algorithm works by having each node maintain a fragment of the MST, and when a link fails, the two nodes connected by that link will merge their fragments to form a new fragment. The new fragment will then be used to find a new MST.

In a network that converged to a MST, when a tree link $e = (u, v)$ fails (with $u$ the
parent of $v$ in the current MST) the spanning tree is split into two \emph{fragments}
that must be reconnected by finding the new minimum-weight outgoing edge (MOE) between
them.

The failure response proceeds in four phases.

\paragraph{Detection and split.}
Both endpoints are notified of the failure.

Both endpoints generate a new unique fragment identity. %(we should choose where to explain this)

Node $u$ (the parent) propagates the information of the failure upward toward the root
$r$ of its fragment. Upon receiving this notification, $r$ enters the \emph{reiden} state.

Node $v$ (the child) declares itself the root of its new fragment and enters the
\emph{reiden} state.

\paragraph{RE-IDENtification.}
The root broadcasts the change down the tree so that every node in the fragment learns the
new identity. When nodes change their identity, they acknowledge the change to their
parent, such that the root can determine when all nodes have been updated.

\paragraph{Minimum outgoing edge search.}
When the root has received acknowledgements from all its children, it starts the search
for the MOE.

The root starts an echo wave down the tree to collect information about outgoing edges.

Each node receiving the echo wave:
\begin{itemize}
\item collects its outgoing edges (by comparing ids).
\item decides its local MOE (the minimum-weight outgoing edge among its own outgoing edges
and those reported by its children).
\item sends the local MOE to its parent.
\end{itemize}

When the root has received all replies, it determines the global MOE of the fragment.

\paragraph{Fragment merge.}
The root sends a message down the tree to inform the node at the MOE that it has been
selected as the fragment's MOE.

Now:
- if no other changes: continue like ghs
- otherwise ...

\section{Physical architecture and deployment}
Which nodes and platforms involved, and where each component is deployed.

\section{Development plan}
Since it is difficult to predict just how hard implementing a new system will be, you should formulate as a set of ``tiers,'' where the basic tier is something you’re sure you can complete, and the additional tiers add more features, at both the application and the system level.

\chapter{Implementation}

Details about the implementation: every choice about platforms, languages, software/hardware, middlewares, which has not been decided in the requirements.


Important choices about implementation should be described here; e.g., peculiar data structures.


\chapter{Validation}

Check if requirements from Chapter~\ref{ch:analysis} have been fulfilled.
Quantitative tests (simulations) and screenshots of the interfaces are put here.


\chapter{Conclusions}

What has been done with respect to what has been promised in Chapters~\ref{ch:intro} and \ref{ch:analysis}, and what is left out.

\appendix

\chapter{Appendix}

In the Appendix you can put code snippets, snapshots, installation instructions, etc.


\chapter*{Evaluation}
Your system will be judged mainly on how it operates as a distributed system. The primary evaluation will be according to whether your system has the following attributes:
\begin{itemize}
\item  It should be an interesting distributed system, making use of some of the algorithms we have covered in class for distributed synchronization, replication, fault tolerance and recovery, security, etc.
\item The software should be well designed and well implemented, in terms of the overall architecture and the detailed realization.
\item You should devise and apply systematic testing procedures, at both the unit and systems levels.
\item The system should operate reliably and with good performance, even in the face of failures.
\end{itemize}
Important, but secondary considerations include:
\begin{itemize}
\item Time taken to do the project (the sooner the better, but do not miss details in order to end sooner)
\item  How nice is the application's appearance: does it have a nice interface or a compelling visual display?
\end{itemize}

\end{document}